\documentclass[10pt,a4paper]{amsart}
\usepackage{amsmath,amssymb,amsthm,graphicx}
% Retained minimal preamble for the Stacks Project chapter more-morphisms.tex (More on Morphisms).
% Verbatim excerpt of the chapter preamble: Xy-pic setup (native preamble lines 14-18) and the native
% theorem declarations with the mathematical macros (native preamble lines 154-217). Copied mechanically;
% nothing added, removed or redefined. The wrapper class and the amsmath/amssymb/amsthm/graphicx base
% packages are supplied by the assembler, so the class/externaldocument part of the native preamble is omitted.
% For commutative diagrams we use Xy-pic
\usepackage[all]{xy}

% We use 2cell for 2-commutative diagrams.
\xyoption{2cell}
% Theorem environments.
%
\theoremstyle{plain}
\newtheorem{theorem}[subsection]{Theorem}
\newtheorem{proposition}[subsection]{Proposition}
\newtheorem{lemma}[subsection]{Lemma}

\theoremstyle{definition}
\newtheorem{definition}[subsection]{Definition}
\newtheorem{example}[subsection]{Example}
\newtheorem{exercise}[subsection]{Exercise}
\newtheorem{situation}[subsection]{Situation}

\theoremstyle{remark}
\newtheorem{remark}[subsection]{Remark}
\newtheorem{remarks}[subsection]{Remarks}

\numberwithin{equation}{subsection}

% Macros
%
\def\lim{\mathop{\mathrm{lim}}\nolimits}
\def\colim{\mathop{\mathrm{colim}}\nolimits}
\def\Spec{\mathop{\mathrm{Spec}}}
\def\Hom{\mathop{\mathrm{Hom}}\nolimits}
\def\Ext{\mathop{\mathrm{Ext}}\nolimits}
\def\SheafHom{\mathop{\mathcal{H}\!\mathit{om}}\nolimits}
\def\SheafExt{\mathop{\mathcal{E}\!\mathit{xt}}\nolimits}
\def\Sch{\mathit{Sch}}
\def\Mor{\mathop{\mathrm{Mor}}\nolimits}
\def\Ob{\mathop{\mathrm{Ob}}\nolimits}
\def\Sh{\mathop{\mathit{Sh}}\nolimits}
\def\NL{\mathop{N\!L}\nolimits}
\def\CH{\mathop{\mathrm{CH}}\nolimits}
\def\proetale{{pro\text{-}\acute{e}tale}}
\def\etale{{\acute{e}tale}}
\def\QCoh{\mathit{QCoh}}
\def\Ker{\mathop{\mathrm{Ker}}}
\def\Im{\mathop{\mathrm{Im}}}
\def\Coker{\mathop{\mathrm{Coker}}}
\def\Coim{\mathop{\mathrm{Coim}}}

% Boxtimes
%
\DeclareMathSymbol{\boxtimes}{\mathbin}{AMSa}{"02}

%
% Macros for moduli stacks/spaces
%
\def\QCohstack{\mathcal{QC}\!\mathit{oh}}
\def\Cohstack{\mathcal{C}\!\mathit{oh}}
\def\Spacesstack{\mathcal{S}\!\mathit{paces}}
\def\Quotfunctor{\mathrm{Quot}}
\def\Hilbfunctor{\mathrm{Hilb}}
\def\Curvesstack{\mathcal{C}\!\mathit{urves}}
\def\Polarizedstack{\mathcal{P}\!\mathit{olarized}}
\def\Complexesstack{\mathcal{C}\!\mathit{omplexes}}
% \Pic is the operator that assigns to X its picard group, usage \Pic(X)
% \Picardstack_{X/B} denotes the Picard stack of X over B
% \Picardfunctor_{X/B} denotes the Picard functor of X over B
\def\Pic{\mathop{\mathrm{Pic}}\nolimits}
\def\Picardstack{\mathcal{P}\!\mathit{ic}}
\def\Picardfunctor{\mathrm{Pic}}
\def\Deformationcategory{\mathcal{D}\!\mathit{ef}}
\usepackage{cleveref}
\crefname{originalreference}{original reference}{original references}
\Crefname{originalreference}{Original reference}{Original references}
\makeatletter
\newlabel{cohomology-lemma-tor-amplitude-stalk}{{09U9}{source}{}{sourcepointer-1}{}}
\newlabel{cohomology-lemma-tor-amplitude-stalk@cref}{{[originalreference][0][]09U9}{[1][0][]source}{}{}{}}
\makeatother
\crefname{originalreference}{original reference}{original references}
\Crefname{originalreference}{Original reference}{Original references}
\makeatletter
\newlabel{more-algebra-lemma-check-perfect-stalks}{{068W}{source}{}{sourcepointer-2}{}}
\newlabel{more-algebra-lemma-check-perfect-stalks@cref}{{[originalreference][0][]068W}{[1][0][]source}{}{}{}}
\makeatother
\crefname{originalreference}{original reference}{original references}
\Crefname{originalreference}{Original reference}{Original references}
\makeatletter
\newlabel{more-algebra-lemma-perfect-over-polynomial-ring}{{068X}{source}{}{sourcepointer-3}{}}
\newlabel{more-algebra-lemma-perfect-over-polynomial-ring@cref}{{[originalreference][0][]068X}{[1][0][]source}{}{}{}}
\makeatother
\crefname{originalreference}{original reference}{original references}
\Crefname{originalreference}{Original reference}{Original references}
\makeatletter
\newlabel{more-algebra-lemma-tor-amplitude-localization}{{0B67}{source}{}{sourcepointer-4}{}}
\newlabel{more-algebra-lemma-tor-amplitude-localization@cref}{{[originalreference][0][]0B67}{[1][0][]source}{}{}{}}
\makeatother
\title{Perfectness is checked on the stalks of a pseudo-coherent morphism}
\author{The Stacks Project Authors}
\date{stacks-project:more-morphisms.tex:local-snapshot-20261011; GNU FDL 1.2 or later}
\begin{document}
\maketitle
\noindent\textit{Editorial note.} Counted span is the author's own proof environment, source lines 17913-17951, immediately following the statement at 17900-17911.\par\medskip
\noindent\textit{Editorial note.} Four cross-chapter pointers (Cohomology, More on Algebra) are recorded in assembly.external\_references.\par\medskip
\noindent\textit{Editorial note.} The reduction to the affine case and the polynomial ring representation are displayed inside the span.\par\medskip
\noindent\textit{Allowed standard prerequisites (editorial).}\begin{itemize}
\item Tor amplitude on stalks. Cohomology, lemma cohomology-lemma-tor-amplitude-stalk, cited in the proof. Converts a local tor dimension condition into a stalkwise one.
\item Localisation and tor amplitude. More on Algebra, lemma more-algebra-lemma-tor-amplitude-localization, cited in the proof. Moves tor amplitude along localisations of the ring map.
\item Perfectness over a polynomial ring. More on Algebra, lemma more-algebra-lemma-perfect-over-polynomial-ring, cited at proof lines 17943-17946. Upgrades finite tor dimension over the localised polynomial ring to perfectness there.
\item Checking perfectness on stalks of a finitely presented algebra. More on Algebra, lemma more-algebra-lemma-check-perfect-stalks, cited at proof lines 17948-17949. Concludes that the module is perfect, hence the ring map is perfect, from the stalkwise information.
\end{itemize}\medskip
\ifdefined\hypertarget\hypertarget{sourcepointer-1}{}\fi
\noindent\textit{Source reference.} Original reference 09U9: focused-seeds/geometry/S9-morph-check-perfect-stalks/ref-tags.txt, source lines 4--4. Navigation pointer only: the counted proof cites this statement by \textbackslash{}ref as a standing input; the deduction the counted proof has to make is carried out inside the counted span, and no part of that deduction is delegated to this pointer.\par\medskip
\ifdefined\hypertarget\hypertarget{sourcepointer-2}{}\fi
\noindent\textit{Source reference.} Original reference 068W: focused-seeds/geometry/S9-morph-check-perfect-stalks/ref-tags.txt, source lines 5--5. Navigation pointer only: the counted proof cites this statement by \textbackslash{}ref as a standing input; the deduction the counted proof has to make is carried out inside the counted span, and no part of that deduction is delegated to this pointer.\par\medskip
\ifdefined\hypertarget\hypertarget{sourcepointer-3}{}\fi
\noindent\textit{Source reference.} Original reference 068X: focused-seeds/geometry/S9-morph-check-perfect-stalks/ref-tags.txt, source lines 6--6. Navigation pointer only: the counted proof cites this statement by \textbackslash{}ref as a standing input; the deduction the counted proof has to make is carried out inside the counted span, and no part of that deduction is delegated to this pointer.\par\medskip
\ifdefined\hypertarget\hypertarget{sourcepointer-4}{}\fi
\noindent\textit{Source reference.} Original reference 0B67: focused-seeds/geometry/S9-morph-check-perfect-stalks/ref-tags.txt, source lines 7--7. Navigation pointer only: the counted proof cites this statement by \textbackslash{}ref as a standing input; the deduction the counted proof has to make is carried out inside the counted span, and no part of that deduction is delegated to this pointer.\par\medskip
\begin{lemma}
\label{lemma-check-perfect-stalks}
Let $f : X \to S$ be a pseudo-coherent morphism of schemes.
The following are equivalent
\begin{enumerate}
\item $f$ is perfect,
\item $\mathcal{O}_X$ locally has finite tor dimension as a
sheaf of $f^{-1}\mathcal{O}_S$-modules, and
\item for all $x \in X$ the ring $\mathcal{O}_{X, x}$ has finite tor
dimension as an $\mathcal{O}_{S, f(x)}$-module.
\end{enumerate}
\end{lemma}

\begin{proof}
The problem is local on $X$ and $S$. Hence we may assume that
$X = \Spec(B)$, $S = \Spec(A)$ and $f$ corresponds to a pseudo-coherent
ring map $A \to B$.

\medskip\noindent
If (1) holds, then $B$ has finite tor dimension $d$ as $A$-module. Then
$B_\mathfrak q$ has tor dimension $d$ as an $A_\mathfrak p$-module for all
primes $\mathfrak q \subset B$ with $\mathfrak p = A \cap \mathfrak q$, see
More on Algebra, Lemma \ref{more-algebra-lemma-tor-amplitude-localization}.
Then $\mathcal{O}_X$ has tor dimension $d$ as a 
sheaf of $f^{-1}\mathcal{O}_S$-modules by
Cohomology, Lemma \ref{cohomology-lemma-tor-amplitude-stalk}.
Thus (1) implies (2).

\medskip\noindent
By Cohomology, Lemma \ref{cohomology-lemma-tor-amplitude-stalk} (2) implies
(3).

\medskip\noindent
Assume (3). We cannot use
More on Algebra, Lemma \ref{more-algebra-lemma-tor-amplitude-localization}
to conclude as we are not given that the tor dimension of
$B_\mathfrak q$ over $A_\mathfrak p$ is bounded independent of $\mathfrak q$.
Choose a presentation $A[x_1, \ldots, x_n] \to B$. Then $B$ is
pseudo-coherent as a $A[x_1, \ldots, x_n]$-module. Let
$\mathfrak q \subset A[x_1, \ldots, x_n]$ be a prime ideal
lying over $\mathfrak p \subset A$. Then either $B_\mathfrak q$ is zero
or by assumption it has finite tor dimension as an
$A_\mathfrak p$-module. Since the fibres of $A \to A[x_1, \ldots, x_n]$
have finite global dimension, we can apply
More on Algebra, Lemma \ref{more-algebra-lemma-perfect-over-polynomial-ring}
to $A_\mathfrak p \to A[x_1, \ldots, x_n]_\mathfrak q$
to see that $B_\mathfrak q$ is a perfect
$A[x_1, \ldots, x_n]_\mathfrak q$-module. Hence
$B$ is a perfect $A[x_1, \ldots, x_n]$-module by
More on Algebra, Lemma \ref{more-algebra-lemma-check-perfect-stalks}.
Thus $A \to B$ is a perfect ring map by definition.
\end{proof}

\end{document}
